\chapter{Self-Refinement and Scoring}\label{ch:refine}

\tikzset{
  r-stage/.style={draw=sysblue, rounded corners=2pt, fill=sysbluebg, minimum height=8mm,
    text width=2.4cm, align=center, font=\footnotesize},
  r-hot/.style={r-stage, fill=keygreen!22, draw=keygreen, very thick},
  r-arrow/.style={-{Stealth[length=2.2mm]}, thick, color=black!65},
  r-num/.style={circle, fill=labpurple, text=white, font=\scriptsize\bfseries, inner sep=1.3pt},
  r-box/.style={draw=black!45, rounded corners=2pt, fill=white, font=\scriptsize,
    align=left, inner sep=4pt},
}

The two inference-time methods of \cref{ch:scaling} both work \emph{in parallel}: one long answer, or
many independent answers and a vote. Neither lets the model look at what it wrote and fix it. People
solving hard problems do exactly that --- draft, check, revise --- and this chapter gives the model
the same loop. \emph{Self-refinement} asks the model to critique its own draft and then to produce a
revised answer that addresses the critique.

Refinement immediately raises a question that voting avoided: \emph{is the revision actually
better?} Models can talk themselves out of a correct answer as easily as into one. So before building
the loop we build \emph{scorers} --- functions that assign a number to a response --- starting from
a transparent rule-based score and moving to the model's own confidence, measured through token
log-probabilities. Log-probabilities are worth the careful treatment they get here: in
\cref{ch:grpo} they reappear inside the RL training objective, and in \cref{ch:distill} inside the
distillation loss. From the systems side, scoring is the cheap half of inference --- a single
prefill, no decode loop --- and refinement is the expensive, strictly sequential half; we will
quantify both.

\begin{figure}[t]
\centering
\begin{tikzpicture}[node distance=4mm and 4mm]
\node[r-stage] (load) {Load model\\(with logits)};
\node[r-stage, right=of load] (rule) {Rule-based\\score};
\node[r-stage, right=of rule] (prob) {Token\\probabilities};
\node[r-stage, right=of prob] (logp) {Token log\\probabilities};
\node[r-stage, below=8mm of logp] (conf) {Log-prob\\confidence score};
\node[r-hot, left=of conf] (refine) {Self-refinement\\(critique, revise)};
\node[r-hot, left=of refine] (loop) {Score-gated\\refinement loop};
\node[r-stage, left=of loop, fill=fcompute!15] (sys) {Cost, latency,\\search};
\foreach \n/\k in {load/1, rule/2, prob/3, logp/4, conf/5, refine/6, loop/7, sys/8} {\node[r-num] at (\n.north west) {\k};}
\draw[r-arrow] (load) -- (rule); \draw[r-arrow] (rule) -- (prob); \draw[r-arrow] (prob) -- (logp);
\draw[r-arrow] (logp) -- (conf); \draw[r-arrow] (conf) -- (refine); \draw[r-arrow] (refine) -- (loop);
\draw[r-arrow] (loop) -- (sys);
\draw[decorate, decoration={brace, amplitude=4pt}, thick, sysblue]
  ($(rule.north west)+(0,0.35)$) -- ($(logp.north east)+(0,0.35)$)
  node[midway, above=4pt, font=\scriptsize\color{sysblue}] {scorers: how good is a response?};
\end{tikzpicture}
\caption{Roadmap. Stages 2--5 build scoring functions; stages 6--7 build the refinement loop that
uses them to accept or reject revisions; stage 8 is the systems view.}
\label{fig:refine:roadmap}
\end{figure}

% ======================================================================
\section{Scoring and iteratively improving responses}\label{sec:refine:overview}
\covers{RM 5.1}

\Cref{fig:refine:three} places self-refinement next to the two methods we already have. They differ
in the \emph{axis} along which they spend extra compute: chain-of-thought spends it on a longer
single answer, self-consistency on more answers, self-refinement on more rounds of revision.

\begin{figure}[t]
\centering
\begin{tikzpicture}[x=1cm, y=1cm]
% CoT
\node[font=\footnotesize\bfseries] at (1.6,2.6) {Chain of thought};
\node[font=\scriptsize, color=black!60] at (1.6,2.2) {more tokens per answer};
\fill[keygreen!20] (0.0,1.4) rectangle (3.2,1.8); \draw (0.0,1.4) rectangle (3.2,1.8);
\node[font=\tiny] at (1.6,1.6) {long reasoning, one answer};
% SC
\node[font=\footnotesize\bfseries] at (6.0,2.6) {Self-consistency};
\node[font=\scriptsize, color=black!60] at (6.0,2.2) {more answers per question};
\foreach \i in {0,...,3} {
  \fill[sysblue!20] (4.7,{1.8-0.32*\i}) rectangle (6.9,{1.55-0.32*\i});
  \draw (4.7,{1.8-0.32*\i}) rectangle (6.9,{1.55-0.32*\i});
}
\node[draw, circle, inner sep=1pt, font=\tiny] at (7.35,1.2) {vote};
% SR
\node[font=\footnotesize\bfseries] at (10.3,2.6) {Self-refinement};
\node[font=\scriptsize, color=black!60] at (10.3,2.2) {more rounds per answer};
\foreach \i/\t in {0/draft, 1/critique, 2/revise} {
  \fill[napkin!20] ({8.5+1.35*\i},1.4) rectangle ({9.55+1.35*\i},1.8);
  \draw ({8.5+1.35*\i},1.4) rectangle ({9.55+1.35*\i},1.8);
  \node[font=\tiny] at ({9.025+1.35*\i},1.6) {\t};
}
\draw[r-arrow] (9.55,1.6) -- (9.85,1.6); \draw[r-arrow] (10.9,1.6) -- (11.2,1.6);
\draw[r-arrow] (11.725,1.4) -- (11.725,1.05) -- node[below, font=\tiny]{repeat} (10.375,1.05) -- (10.375,1.38);
\draw[thick, black!40] (0,0.45) -- (12,0.45);
\node[font=\scriptsize, align=center] at (1.6,0.1) {latency $\propto$ length};
\node[font=\scriptsize, align=center] at (6.0,0.1) {latency $\approx$ one answer (batched)};
\node[font=\scriptsize, align=center] at (10.3,0.1) {latency $\propto$ rounds};
\end{tikzpicture}
\caption{Three ways to spend inference compute. Self-refinement is the only one that conditions later
work on an explicit evaluation of earlier work; it is also inherently sequential.}
\label{fig:refine:three}
\end{figure}

Self-consistency has become a common choice when accuracy matters more than latency, but it has a
structural limitation: it needs short final answers that can be compared for equality. A scorer
removes that limitation. Given any function $s(\text{question}, \text{response}) \in \R$ we can:
\begin{itemize}
  \item \textbf{rank} candidates and return the best (best-of-$n$, reranking);
  \item \textbf{gate} a revision: accept it only if it scores higher than the draft;
  \item \textbf{guide search}: expand only the most promising partial solutions (\cref{sec:refine:search}).
\end{itemize}
Scorers come in four flavours, in increasing cost: heuristic rules, the model's own token
probabilities, external verifiers (\cref{ch:eval}), and learned reward models or LLM judges.

% ======================================================================
\section{Loading a model whose logits we can see}\label{sec:refine:load}
\covers{RM 5.2}

We load Qwen3-0.6B exactly as in \cref{sec:eval:load}. One requirement is new: the confidence scorers
need the model's full next-token distributions, not just sampled text. With local weights that is
free --- the forward pass returns logits for every position. Hosted APIs typically expose at most the
log-probabilities of the top few tokens per position, and only for generated text, which is enough
for some scorers in this chapter but not all.

\begin{systemslens}[Logit access is a deployment decision]
Returning full logits from an inference server is expensive: $\vocab \times 4$\,B $= 0.6$\,MB per
token for Qwen3's vocabulary, i.e.\ 600\,MB for a 1000-token response --- far more than the text.
Serving engines therefore return \emph{selected} log-probs (of the sampled token and the top-$k$
alternatives) computed on the GPU, and offer a separate ``prompt log-probs'' option that scores a
given text in a prefill-only pass. Designing scorers that need only the log-prob of the realized
tokens, as ours do, keeps them deployable.
\end{systemslens}

% ======================================================================
\section{A rule-based score}\label{sec:refine:rule}
\covers{RM 5.3}

The simplest scorer encodes preferences as rules. For math responses we prefer, in this order, a
response with a final boxed answer, then one that at least contains a number; and among those, shorter
ones. \Cref{lst:refine:heuristic} implements this with a format bonus plus a smooth brevity term
$\beta/(\beta + \ell)$ that lies in $(0, 1]$, so format always dominates length.

\begin{lstlisting}[caption={A rule-based scorer for format and brevity (\texttt{llmsys/tts.py}).},label={lst:refine:heuristic}]
def heuristic_score(answer, boxed_bonus=2.0, number_bonus=1.0,
                    brevity_scale=500.0):
    s = 0.0
    if last_boxed(answer) is not None:          # complete \boxed{...}
        s += boxed_bonus
    elif NUMBER_RE.search(answer or ""):        # at least some number
        s += number_bonus
    return s + brevity_scale / (brevity_scale + len(answer or ""))
\end{lstlisting}

\begin{table}[t]
\centering\small
\begin{tabular}{@{}p{6.4cm}rrr@{}}
\toprule
Response (abbreviated) & format & brevity & score \\
\midrule
``\dots so $ab = 60$. \texttt{\textbackslash boxed\{60\}}'' (180 chars) & 2.0 & 0.735 & 2.735 \\
``Let $a,b$ be the sides \dots{} {[}long{]} \dots{} \texttt{\textbackslash boxed\{60\}}'' (900 chars) & 2.0 & 0.357 & 2.357 \\
``The area is 60.'' (16 chars) & 1.0 & 0.969 & 1.969 \\
``I am not sure how to proceed.'' (29 chars) & 0.0 & 0.945 & 0.945 \\
\bottomrule
\end{tabular}
\caption{Heuristic scores for four responses with $\beta = 500$. The score measures format and
brevity, not correctness: a short boxed \emph{wrong} answer would beat a long correct one.}
\label{tab:refine:heuristic}
\end{table}

\begin{pitfall}[Goodhart's law in one function]
Any score you optimize against will be exploited. Gate refinements on brevity and the model learns to
delete its reasoning; reward boxes and it boxes a guess. Rule-based scores are useful as cheap
filters, tie-breakers and format checks (\cref{ch:grpo2} adds exactly such a format reward to RL),
but they should never be the only signal of quality.
\end{pitfall}

% ======================================================================
\section{Token probabilities}\label{sec:refine:tokprob}
\covers{RM 5.4}

A more informative signal is the model's own \emph{confidence}: how probable does the model find
the tokens of a given answer? Recall from \cref{sec:scaling:select} that at every position the model
produces logits over the whole vocabulary, which the softmax turns into a next-token distribution.
To generate, we pick or sample from that distribution. To \emph{score}, we feed in a complete
candidate response and simply look up the probability the model assigned to each token that is
actually there --- \emph{teacher forcing}. No sampling happens; we are only inspecting the
distribution.

Consider two candidate completions of the prompt ``The capital of France is'': `` Paris.'' and
`` Pasta.''. A well-trained model assigns most of its next-token probability after ``is'' to
`` Paris'' and almost none to `` Pasta''. By the chain rule of probability, the probability of a
whole answer $y = (y_1, \dots, y_m)$ given a prompt $x$ is the product of per-token conditionals,
\begin{equation}
  \Pr_\theta(y \mid x) \;=\; \prod_{t=1}^{m} \Pr_\theta\bigl(y_t \mid x, y_{<t}\bigr),
  \label{eq:refine:chain}
\end{equation}
and each factor is read off one row of the logits tensor (\cref{fig:refine:shift}).

\begin{figure}[t]
\centering
\begin{tikzpicture}[tok/.style={draw=black!40, minimum width=1.25cm, minimum height=6mm,
  font=\scriptsize\ttfamily, fill=white, inner sep=1pt}]
\foreach \t/\x/\c in {The/0/black!6, { capital}/1.35/black!6, { of}/2.7/black!6,
  { France}/4.05/black!6, { is}/5.4/black!6, { Paris}/6.75/keygreen!20, {.}/8.1/keygreen!20} {
  \node[tok, fill=\c] at (\x,0) {\t};
}
\node[font=\scriptsize, anchor=east] at (-0.75,0) {input ids};
\foreach \x/\l in {0/{p( capital)}, 1.35/{p( of)}, 2.7/{p( France)}, 4.05/{p( is)}, 5.4/{p( Paris)}, 6.75/{p(.)}} {
  \node[tok, fill=sysbluebg, font=\tiny\ttfamily] (r\x) at (\x,-1.3) {\l};
  \draw[r-arrow, black!35] (\x,-0.3) -- (\x,-1.0);
}
\node[tok, fill=black!5, font=\tiny] at (8.1,-1.3) {(unused)};
\node[font=\scriptsize, anchor=east] at (-0.75,-1.3) {logits row $i$};
\draw[decorate, decoration={brace, amplitude=4pt, mirror}, thick, keygreen]
  (4.75,-1.7) -- (7.4,-1.7) node[midway, below=4pt, font=\scriptsize\color{keygreen}] {answer tokens};
\draw[decorate, decoration={brace, amplitude=4pt, mirror}, thick, black!50]
  (-0.6,-1.7) -- (4.65,-1.7) node[midway, below=4pt, font=\scriptsize\color{black!60}] {prompt positions (excluded from the score)};
\end{tikzpicture}
\caption{Teacher forcing: row $i$ of the logits predicts token $i{+}1$. One forward pass over prompt
plus answer yields the probability of every answer token; the prompt's own predictions are discarded
and the last row predicts beyond the sequence.}
\label{fig:refine:shift}
\end{figure}

The off-by-one in \cref{fig:refine:shift} is the single most common bug in log-probability code, in
scoring and in RL alike. The rule: logits at position $i$ are the distribution for the token at
position $i+1$, so we drop the last logits row and the first token, and the answer's first token is
predicted by the \emph{last prompt position}.

\begin{lstlisting}[caption={Per-token probabilities of an answer by teacher forcing.},label={lst:refine:tokprob}]
@torch.no_grad()
def answer_token_probs(model, prompt_ids, answer_ids):
    ids = torch.tensor([list(prompt_ids) + list(answer_ids)])
    logits = model(ids)[0, :-1].float()            # row i predicts token i+1
    probs = torch.softmax(logits, dim=-1)
    targets = ids[0, 1:]                           # the tokens actually present
    tok_p = probs.gather(-1, targets[:, None]).squeeze(-1)
    return tok_p[len(prompt_ids) - 1:]             # answer tokens only
\end{lstlisting}

\begin{keyidea}[Confidence is not correctness]
A high probability means the answer fits the patterns the model learned, not that it is true. Models
can be confidently wrong (a common misconception repeated in training data) and hesitantly right
(an unusual but correct phrasing). Probability-based scores are useful because, on average and for
well-trained models, confidence correlates with correctness; they are a weak verifier, not a
verifier.
\end{keyidea}

\begin{systemslens}[Scoring is prefill, not generation]
Generating an $m$-token answer takes $m$ sequential decode steps, each memory-bound. Scoring the same
answer takes \emph{one} forward pass over $|x| + m$ tokens: all positions are processed in parallel
with large matmuls, which is compute-bound and runs near peak FLOP/s. On an H100, a 1000-token
scoring pass for Qwen3-0.6B is about $2 \times 0.6\text{B} \times 1000 = 1.2$\,TFLOP, roughly a
millisecond of tensor-core time at a realistic fraction of peak, whereas generating those 1000 tokens
takes at least 1000 memory-bound steps ($\gtrsim 0.4$\,s). Scoring is two to three orders of magnitude
cheaper than generating, which is why reward models, rerankers and PRMs are deployed as prefill-only
services that batch many candidates together.
\end{systemslens}

% ======================================================================
\section{From probabilities to log-probabilities}\label{sec:refine:logprob}
\covers{RM 5.5}

Multiplying hundreds of probabilities in \cref{eq:refine:chain} quickly produces numbers too small to
represent. We work with logarithms instead: the log of a product is a sum,
\begin{equation}
  \log \Pr_\theta(y \mid x) \;=\; \sum_{t=1}^{m} \log \Pr_\theta(y_t \mid x, y_{<t}),
\end{equation}
and sums of moderately negative numbers are perfectly well behaved. Because $\log$ is strictly
increasing, it preserves every ordering: a more probable sequence always has a larger (less
negative) log-probability. A log-probability of $0$ means certainty; more negative means less likely
(\cref{fig:refine:three-views}).

\begin{figure}[t]
\centering
\begin{tikzpicture}
\begin{groupplot}[group style={group size=3 by 1, horizontal sep=1.3cm},
  width=0.36\linewidth, height=4.2cm, xtick={1,...,7}, xticklabels={a,b,c,d,e,f,g},
  x tick label style={font=\scriptsize\ttfamily}, title style={font=\footnotesize},
  tick label style={font=\scriptsize}]
\nextgroupplot[ybar, bar width=6pt, title={logits $z_i$}, ymin=-1.2, ymax=2.2]
\addplot[fill=sysblue!55, draw=sysblue] coordinates {(1,-1) (2,-0.5) (3,0) (4,0.5) (5,1) (6,1.5) (7,2)};
\nextgroupplot[ybar, bar width=6pt, title={$\softmax(z)_i$}, ymin=0, ymax=0.5]
\addplot[fill=keygreen!55, draw=keygreen] coordinates {(1,0.02) (2,0.033) (3,0.055) (4,0.091) (5,0.149) (6,0.246) (7,0.406)};
\nextgroupplot[ybar, bar width=6pt, title={$\log\softmax(z)_i$}, ymin=-4.2, ymax=0]
\addplot[fill=napkin!55, draw=napkin] coordinates {(1,-3.902) (2,-3.402) (3,-2.902) (4,-2.402) (5,-1.902) (6,-1.402) (7,-0.902)};
\end{groupplot}
\end{tikzpicture}
\caption{Seven logits, their probabilities and their log-probabilities (computed for this figure).
All three preserve the same ordering; log-softmax is simply the logits shifted by a constant,
$\log\softmax(z)_i = z_i - \log\sum_j e^{z_j}$.}
\label{fig:refine:three-views}
\end{figure}

\begin{napkinmath}[When does the product underflow?]
An answer of 1000 tokens with an average per-token probability of $0.9$ has joint probability
$0.9^{1000} = e^{-105.4} \approx 1.7\times10^{-46}$ --- below float32's smallest normal number
($1.2\times10^{-38}$), so it underflows to a subnormal and then to zero. In float16 (smallest normal
$6.1\times10^{-5}$) the product is already subnormal after about 90 such tokens. The log-probability,
$-105.4$, is an ordinary number in any format. Rule: never multiply probabilities; sum
log-probabilities.
\end{napkinmath}

Computing \texttt{torch.log(torch.softmax(z))} is numerically poor: a probability that underflows to
zero becomes $-\infty$ after the log. \texttt{torch.log\_softmax} uses the identity in the caption of
\cref{fig:refine:three-views} together with the log-sum-exp trick,
\begin{equation}
  \log \sum_j e^{z_j} \;=\; M + \log \sum_j e^{z_j - M}, \qquad M = \max_j z_j,
  \label{eq:refine:lse}
\end{equation}
so every exponent is $\le 0$ (no overflow) and the largest term is exactly $1$ (no catastrophic
underflow of the sum). It is also faster: one fused kernel instead of two passes.

\begin{lstlisting}[caption={Per-token log-probabilities in one pass (\texttt{llmsys/tts.py}).},label={lst:refine:tokenlp}]
@torch.no_grad()
def token_logprobs(model, prompt_ids, answer_ids):
    ids = torch.tensor([list(prompt_ids) + list(answer_ids)])
    logits = model(ids)[0, :-1].float()             # upcast BEFORE log_softmax
    lp = torch.log_softmax(logits, -1).gather(-1, ids[0, 1:, None]).squeeze(-1)
    return lp[len(prompt_ids) - 1:]
\end{lstlisting}

\begin{systemslens}[Never materialize the full log-softmax]
\texttt{log\_softmax} over the full $T \times \vocab$ logits creates a float32 tensor of
$T \times \num{151936} \times 4$\,B --- 1.2\,GB for a 2000-token sequence --- only to read one element
per row. In training (\cref{ch:grpo}) this tensor is also kept for the backward pass, and for a batch
it easily exceeds the memory of the model itself. Efficient implementations compute
$z_{y_t} - \mathrm{logsumexp}(z)$ directly, chunking over positions or vocabulary, or fuse the final
projection with the cross-entropy so full logits are never materialized (``fused linear
cross-entropy'' kernels).
\end{systemslens}

% ======================================================================
\section{Scoring confidence with log-probabilities}\label{sec:refine:confidence}
\covers{RM 5.6}

Two adjustments turn the joint log-probability into a usable confidence score.

\paragraph{Score only the answer.} The prompt is shared by all candidates for a question, so its
tokens add the same constant to each joint log-probability --- but only if they are scored in the
same context. Excluding them, as in \cref{lst:refine:tokenlp}, makes scores comparable and cheaper
to reason about.

\paragraph{Normalize by length.} Every token contributes a non-positive term, so the sum
systematically prefers \emph{shorter} answers. The \emph{average} log-probability per answer token,
\begin{equation}
  \bar\ell(y \mid x) \;=\; \frac{1}{m}\sum_{t=1}^{m} \log \Pr_\theta(y_t \mid x, y_{<t}),
  \qquad \mathrm{PPL}(y \mid x) = e^{-\bar\ell},
  \label{eq:refine:avg}
\end{equation}
is a length-normalized confidence; its exponentiated negative is the familiar \emph{perplexity}
(\cref{ch:pretrain}) of the answer under the model. \Cref{fig:refine:pertoken} shows why the average
is informative: most tokens of a fluent answer are near-certain, and the score is driven by a few
``decision'' tokens --- often exactly the digits of the final result.

\begin{figure}[t]
\centering
\begin{tikzpicture}
\begin{axis}[ybar, width=0.85\linewidth, height=4.8cm, bar width=14pt, ymin=-3, ymax=0.3,
  xtick={1,...,7}, xticklabels={The, {\_area}, {\_is}, {\_}, 6, 0, .}, xmin=0.4, xmax=7.6,
  x tick label style={font=\scriptsize\ttfamily}, ylabel={$\log p(y_t \mid \cdot)$}]
\addplot[fill=sysblue!50, draw=sysblue] coordinates {(1,-0.9) (2,-0.05) (3,-0.01) (4,-0.3) (5,-2.6) (6,-0.02) (7,-0.4)};
\draw[dashed, warnred, thick] (axis cs:0.4,-0.611) -- (axis cs:7.6,-0.611);
\node[font=\scriptsize\color{warnred}, anchor=north east] at (axis cs:7.5,-0.68) {mean $-0.61$};
\node[font=\scriptsize, anchor=north] at (axis cs:5,-2.62) {$-2.6$};
\end{axis}
\end{tikzpicture}
\caption{Illustrative per-token log-probabilities for the answer ``The area is 60.'' The sum is
$-4.28$ (joint probability $0.014$), the mean $-0.61$ (perplexity $1.84$). The uncertain token is the
first digit of the result.}
\label{fig:refine:pertoken}
\end{figure}

\begin{lstlisting}[caption={Confidence summaries from token log-probabilities (Lab 16, Task 3).},label={lst:refine:conf}]
def confidence_scores(token_lps):
    mean = float(token_lps.mean())
    return dict(sum=float(token_lps.sum()),   # joint log-prob (favours short)
                mean=mean,                    # length-normalized confidence
                ppl=math.exp(-mean),          # perplexity of the answer
                min=float(token_lps.min()))   # weakest link

def avg_logprob_score(model, prompt_ids, answer_ids):
    lp = token_logprobs(model, prompt_ids, answer_ids)
    return float(lp.mean()) if lp.numel() else -math.inf
\end{lstlisting}

Which summary to use is an empirical question. The mean is the standard choice. The minimum token
log-probability flags a single shaky step in an otherwise fluent derivation. Restricting the average
to the final-answer tokens (the content of the box) measures confidence in the result rather than in
the prose. Whatever you choose, check it on a validation set: plot accuracy against score deciles and
see whether higher scores really mean more correct answers.

\begin{pitfall}[Scores from different models or prompts are not comparable]
A log-probability is relative to a model \emph{and} a context. Scoring a revision under the
refinement prompt (which contains the critique) and the draft under the original prompt compares
two different conditional distributions. Always score every candidate under the \emph{same} model
and the \emph{same} scoring prompt --- typically the original question, as \texttt{run.py} for Lab~16
does.
\end{pitfall}


\subsection{A worked comparison}
Suppose a base model produced two candidate solutions to the rectangle problem: $A$ ends in
\texttt{\textbackslash boxed\{60\}} after a five-line derivation, and $B$ ends in
\texttt{\textbackslash boxed\{72.25\}} after wrongly assuming a square. Scoring both under the original
question prompt gives each a vector of per-token log-probabilities. The \emph{sum} is dominated by
length: the longer, correct $A$ may well have the lower joint log-probability. The \emph{mean} asks a
fairer question --- which answer did the model find more natural per token? --- and on problems the
model ``knows'', correct derivations tend to have fewer low-probability tokens than flawed ones,
whose errors often show up as a few surprising tokens (the unjustified step, the wrong digit). The
\emph{minimum} isolates exactly those tokens. None of the three is guaranteed to rank $A$ above $B$;
Exercise~\ref{ex:refine:calib} asks you to measure how often each does.

% ======================================================================
\section{Learned scorers: reward models}\label{sec:refine:rm}

Between hand-written rules and the model's raw confidence sits a third option that the next part of
the book relies on: a \emph{learned} scorer. A \textbf{reward model} (RM) is usually a copy of a
language model whose unembedding layer is replaced by a scalar head; it reads (question, response)
and outputs a real number. It is trained on human (or AI) preference pairs $(y^+, y^-)$ with the
Bradley--Terry likelihood from \cref{sec:eval:leaderboards},
\begin{equation}
  \mathcal{L}_{\text{RM}} \;=\; -\log \sigma\bigl(r_\phi(x, y^+) - r_\phi(x, y^-)\bigr),
\end{equation}
so that differences of reward are log-odds of preference. Variants include \emph{outcome} reward
models (one score per full response), \emph{process} reward models (one score per reasoning step,
\cref{sec:eval:beyond}) and \emph{generative} verifiers that write a critique and then a verdict ---
an LLM judge trained for the purpose.

Reward models are the scorer of choice for open-ended tasks and the centrepiece of classic RLHF
(\cref{ch:grpo}). Their weakness is the one every learned function shares: off the training
distribution they are wrong in ways an optimizer will find. Best-of-$n$ with a large $n$, or RL
against an RM for many steps, eventually produces responses that score highly and are bad
(``reward hacking''). Verifiable rewards avoid this for math and code, which is a large part of why
reasoning-model training moved to them.

\begin{systemslens}[Serving a reward model]
An RM forward pass is a prefill with a scalar output: no KV cache needs to be kept, no decode loop
runs, and batches can be packed tightly by total tokens. In RL systems (\cref{ch:rlsys}) the RM is
typically a separate service colocated with or disaggregated from the policy, sized by the number of
rollouts per step times their length. For a 7B RM scoring 1024 rollouts of 1000 tokens, that is
$2 \times 7\times10^9 \times 1.024\times10^6 \approx 1.4\times10^{16}$ FLOPs, about 30\,s of one
H100 at a realistic 50\% of peak --- small next to generating those rollouts, but not free.
\end{systemslens}

% ======================================================================
\section{Self-refinement through iterative feedback}\label{sec:refine:selfrefine}
\covers{RM 5.7}

Self-refinement (Madaan et al., 2023, ``Self-Refine'') uses one model in three roles, each a plain
generation call with a different prompt (\cref{fig:refine:loop}):
\begin{enumerate}
  \item \textbf{Draft.} Answer the question as usual.
  \item \textbf{Critique.} Given the question and the draft, list errors and a short fix plan.
  \item \textbf{Revise.} Given the question, the draft and the critique, write an improved answer
  that ends with a boxed result.
\end{enumerate}
Steps 2--3 can be repeated, each revision becoming the next draft. Nothing here is new machinery ---
it is sequential composition of the generation function from \cref{ch:eval} with prompt templates.

\begin{figure}[t]
\centering
\begin{tikzpicture}[node distance=6mm and 6mm]
\node[r-box, text width=2.6cm, fill=black!4] (q) {\textbf{Question}\\A rectangle has perimeter 34 and diagonal 13. Find its area.};
\node[r-box, text width=2.6cm, right=of q, fill=warnredbg] (d) {\textbf{Draft}\\$a{+}b{=}17$, so $ab = 17^2/4$ \dots\\\texttt{\textbackslash boxed\{72.25\}}};
\node[r-box, text width=2.6cm, right=of d, fill=napkinbg] (c) {\textbf{Critique}\\Assumed a square; the diagonal was never used. Use $a^2{+}b^2{=}169$.};
\node[r-box, text width=2.6cm, right=of c, fill=keygreenbg] (r) {\textbf{Revision}\\$(a{+}b)^2-(a^2{+}b^2)$ $=2ab=120$\\\texttt{\textbackslash boxed\{60\}}};
\draw[r-arrow] (q) -- node[above, font=\tiny]{1} (d);
\draw[r-arrow] (d) -- node[above, font=\tiny]{2} (c);
\draw[r-arrow] (c) -- node[above, font=\tiny]{3} (r);
\node[r-stage, below=9mm of c, text width=3.0cm, fill=fcompute!15] (s) {score draft and revision\\under the same prompt};
\node[r-stage, right=8mm of s, text width=2.6cm, fill=keygreen!20, draw=keygreen] (acc) {keep revision iff\\$s(\text{rev}) > s(\text{draft})$};
\draw[r-arrow] (d.south) |- (s.west);
\draw[r-arrow] (r.south) -- (s.north east);
\draw[r-arrow] (s) -- (acc);
\draw[r-arrow, dashed] (acc.south) -- ++(0,-0.45) -| node[pos=0.25, below, font=\tiny]{winner becomes the next draft (repeat 2--3)} (d.south east);
\end{tikzpicture}
\caption{One round of self-refinement with score gating (example text written for this book). The
draft, critique and revision are three generation calls; the scorer decides which answer survives
into the next round.}
\label{fig:refine:loop}
\end{figure}

\begin{lstlisting}[caption={Critique and refine templates (\texttt{llmsys/tts.py}).},label={lst:refine:templates}]
CRITIQUE_TEMPLATE = (
    "You are reviewing a draft solution.\n\nProblem:\n{question}\n\n"
    "Draft:\n{draft}\n\nList any mistakes or gaps and a short plan to fix them "
    "(at most ~100 words).\n\nReview:")
REFINE_TEMPLATE = (
    "Improve the draft solution using the review. End with the final result as "
    "\\boxed{{ANSWER}}.\n\nProblem:\n{question}\n\nDraft:\n{draft}\n\n"
    "Review:\n{critique}\n\nImproved solution:")
\end{lstlisting}

Walking through one round by hand is the best way to build intuition. Sample a draft at moderate
temperature, print it, generate the critique with a small token budget (a critique that rambles is a
critique that will be ignored), print it, then generate the revision. You will see three patterns:
the critique identifies a real error and the revision fixes it; the critique hallucinates an error
in a correct draft and the revision breaks it; or the critique is generic (``check your arithmetic'')
and the revision is a paraphrase. A scorer exists to filter the second and third cases.

\begin{keyidea}[Feedback quality is the whole game]
Refinement can only be as good as its feedback. With \emph{intrinsic} feedback --- the same model
judging itself, with no external signal --- gains on reasoning tasks are small and sometimes negative:
Huang et al.\ (2024) found that LLMs often fail to self-correct their reasoning without external
feedback, and can even degrade correct answers. Refinement shines when the critique has access to
information the draft lacked: an executed test case, a verifier's verdict, a compiler error, a
retrieved document. In system design terms, self-refinement is a loop; close it with the best
signal you can afford.
\end{keyidea}

% ======================================================================
\section{Coding the refinement loop}\label{sec:refine:loopcode}
\covers{RM 5.8}

\Cref{lst:refine:loop} packages the procedure into a function with an iteration count and an optional
scorer. Without a scorer every revision is accepted. With one, a revision replaces the current answer
only if it scores strictly higher; otherwise the loop critiques the \emph{same} answer again in the
next round (with sampling, the next critique will differ).

\begin{lstlisting}[caption={Score-gated self-refinement (condensed from \texttt{llmsys/tts.py}).},label={lst:refine:loop}]
def self_refine(question, generate_fn, iterations=2, score_fn=None,
                answer_prompt_fn=lambda q: q, max_answer_tokens=1024,
                max_critique_tokens=256):
    cur = generate_fn(answer_prompt_fn(question), max_answer_tokens)
    cur_score = score_fn(question, cur) if score_fn else None
    steps = []
    for it in range(iterations):
        crit = generate_fn(CRITIQUE_TEMPLATE.format(question=question, draft=cur),
                           max_critique_tokens)
        rev = generate_fn(REFINE_TEMPLATE.format(question=question, draft=cur,
                                                 critique=crit), max_answer_tokens)
        rev_score = score_fn(question, rev) if score_fn else None
        ok = score_fn is None or rev_score > cur_score
        steps.append(RefineStep(it, cur, crit, rev, cur_score, rev_score, ok))
        if ok:
            cur, cur_score = rev, rev_score
    return RefineResult(final=cur, final_score=cur_score, steps=steps)
\end{lstlisting}

The scorer for math combines naturally with the pieces of this chapter: score each candidate by its
average token log-probability \emph{under the original question prompt}:

\begin{lstlisting}[caption={Plugging the log-prob scorer into the loop.},label={lst:refine:plug}]
def logprob_scorer(question, answer):
    prompt_ids = tok.encode(render_math_prompt(question))
    return avg_logprob_score(model, prompt_ids, tok.encode(answer))

result = self_refine(problem, generate_fn, iterations=2, score_fn=logprob_scorer,
                     answer_prompt_fn=render_math_prompt)
for s in result.steps:
    print(s.iteration, s.score_draft, s.score_revised, s.accepted)
\end{lstlisting}


\subsection{Acceptance and stopping policies}
Strict improvement is only one of several reasonable acceptance rules, and each encodes a different
belief about the scorer:

\begin{itemize}
  \item \textbf{Always accept} (no scorer): trusts the critique; cheapest; vulnerable to harmful
  revisions.
  \item \textbf{Strict improvement} ($s_{\text{rev}} > s_{\text{cur}}$): trusts the scorer's ranking;
  ties keep the incumbent, which makes the loop conservative.
  \item \textbf{Margin} ($s_{\text{rev}} > s_{\text{cur}} + \delta$): accounts for scorer noise; with
  sampled critiques, a small $\delta$ avoids flip-flopping between near-equivalent answers.
  \item \textbf{Keep the best of all rounds}: score every candidate and return the argmax at the end,
  not the last accepted one. With a fixed scorer this equals strict acceptance, but it also lets you
  combine rounds with best-of-$n$ sampling.
  \item \textbf{Verifier override}: if any candidate passes an exact check (a substitution test, a
  unit test), return it immediately --- an external signal always beats a confidence score.
\end{itemize}

Stopping rules matter as much for cost. Stop when a revision is rejected twice in a row, when the
revised final answer equals the current one (the critique changed only the prose), or when a token
budget is exhausted. Each of these is a one-line change to \cref{lst:refine:loop}, and each can
remove a large fraction of the sequential calls on problems where the first draft was already
right.

\begin{table}[t]
\centering\small
\begin{tabular}{@{}lccc@{}}
\toprule
Policy & needs scorer & protects correct drafts & extra cost \\
\midrule
always accept & no & no & none \\
strict improvement & yes & partly (if scorer agrees) & $2k$ scoring prefills \\
margin $\delta$ & yes & more & same \\
best of all rounds & yes & as strict & same \\
verifier override & exact check & yes, when check applies & check per candidate \\
\bottomrule
\end{tabular}
\caption{Acceptance policies for self-refinement.}
\label{tab:refine:policies}
\end{table}

\begin{pitfall}[What score gating can and cannot fix]
Gating prevents a revision from replacing a draft that the scorer prefers, so it limits the damage
of bad critiques. It cannot fix a scorer that prefers wrong answers. In particular, average
log-probability tends to prefer answers the model finds \emph{familiar}; if the base model's
confident answer is wrong, gating will dutifully keep it. Whether refinement with a given scorer
helps a given model is an empirical question --- measure it on a validation set, including the
fraction of accepted revisions that turned a correct draft into a wrong one.
\end{pitfall}

Refinement composes with the earlier methods. CoT prompting can be used for the draft and the
revision; with a base model it should be. Self-consistency can be applied \emph{after} refinement
(refine $n$ samples, then vote), or refinement can be applied to the vote's winner. Each composition
multiplies cost, which brings us to the systems view.

% ======================================================================
\section{The cost of thinking sequentially}\label{sec:refine:cost}

Let the question have $q$ tokens, every answer $a$ tokens and every critique $c$ tokens, and assume
template text is negligible. Without any KV reuse between calls, one round of refinement prefills
the critique prompt ($q + a$ tokens) and the refine prompt ($q + a + c$), and decodes $c + a$ new
tokens. For $k$ rounds after the draft:
\begin{align}
  \text{prefill tokens} &= q + k\,(2q + 2a + c), \label{eq:refine:prefill}\\
  \text{decode tokens}  &= a + k\,(a + c), \\
  \text{sequential generation calls} &= 1 + 2k.
\end{align}
Scoring adds $2k$ prefill-only passes of $q + a$ tokens if the draft's score is cached.

\begin{napkinmath}[Refinement vs. voting at a matched token budget]
Take $q = 150$, $a = 500$, $c = 120$ and $k = 2$ rounds. Refinement decodes $500 + 2(620) = 1740$
tokens in 5 \emph{sequential} calls and prefills $150 + 2(1420) = 2990$ tokens. Self-consistency with
$n = 4$ decodes $2000$ tokens --- similar --- but in \emph{one} batched call after a single 150-token
prefill. On Qwen3-0.6B at a memory-bound $\approx$0.4\,ms per step, refinement's latency is about
$1740 \times 0.4\,\text{ms} \approx 0.7$\,s of pure decode plus five prefills and Python overhead;
voting's is about $500 \times 0.45\,\text{ms} \approx 0.23$\,s, because its four rows share each
step. Same tokens, roughly $3\times$ the latency: sequential methods buy their ability to condition on
feedback with latency.
\end{napkinmath}

\begin{figure}[t]
\centering
\begin{tikzpicture}[x=0.95cm, y=1cm]
\draw[-{Stealth}] (0,0) -- (12.5,0) node[below left, font=\scriptsize]{time};
% refinement
\node[font=\scriptsize, anchor=east] at (-0.1,1.4) {refinement};
\foreach \a/\b/\c/\l in {0/2.9/sysblue/draft, 2.9/3.6/napkin/crit., 3.6/6.5/keygreen/revise, 6.5/7.2/napkin/crit., 7.2/10.1/keygreen/revise} {
  \fill[\c!30] (\a,1.15) rectangle (\b,1.65); \draw (\a,1.15) rectangle (\b,1.65);
  \node[font=\tiny] at ({(\a+\b)/2},1.4) {\l};
}
\foreach \x in {2.9,6.5,10.1} {\fill[warnred] (\x,1.1) rectangle ({\x+0.12},1.7);}
\node[font=\tiny, color=warnred, anchor=west] at (10.3,1.4) {= scoring (prefill)};
% voting
\node[font=\scriptsize, anchor=east] at (-0.1,0.55) {4-way vote};
\fill[sysblue!30] (0,0.3) rectangle (3.3,0.8); \draw (0,0.3) rectangle (3.3,0.8);
\node[font=\tiny] at (1.65,0.55) {4 samples, one batch};
\end{tikzpicture}
\caption{Schematic timelines at a similar decode-token budget. Refinement's calls cannot overlap
because each consumes the previous output; a batched vote finishes in roughly the time of one
answer.}
\label{fig:refine:timeline}
\end{figure}

\begin{systemslens}[Make the loop cache-friendly]
Every refinement call re-sends text the model has already processed: the question, then the draft.
A serving engine with prefix caching (\cref{ch:serving}) can skip recomputing a prompt prefix it has
seen --- but only an \emph{exact token prefix}. Our templates begin with different instructions
(``You are reviewing\dots'' versus ``Improve the draft\dots''), so the question and draft never form a
shared prefix and every call pays full prefill. Reordering templates so that the stable content comes
first (\texttt{Problem: \dots\ Draft: \dots} followed by the role-specific instruction) turns most of
each refinement prefill into a cache hit. Better still, run refinement as a multi-turn conversation
in which each step appends to the same context: then the entire history is a prefix of the next call
and only new tokens are prefilled. This is a general rule for agentic and multi-call LLM programs:
put the invariant parts of a prompt first.
\end{systemslens}

% ======================================================================
\section{Beyond a single chain: search}\label{sec:refine:search}

Refinement improves one answer along one path. \emph{Search} methods keep several partial solutions
and use a scorer to decide which to extend:

\begin{itemize}
  \item \textbf{Best-of-$n$ reranking}: sample $n$ full answers (in parallel, \cref{sec:scaling:prefix})
  and return the highest-scoring one. With an exact verifier it is unbeatable; with a learned reward
  model it is vulnerable to reward hacking as $n$ grows (the argmax finds the scorer's errors).
  \item \textbf{Step-level beam search}: split solutions into steps (e.g.\ lines), sample several
  continuations per step, keep the top-$b$ partial solutions according to a process reward model
  (PRM) that scores prefixes.
  \item \textbf{Tree search}: Tree-of-Thoughts (Yao et al., 2023) and MCTS-style methods explore a
  tree of partial reasoning states with backtracking, trading many more model calls for the ability
  to recover from early mistakes.
\end{itemize}

All of these are compositions of two primitives this chapter and the last made efficient:
\emph{parallel sampling from a shared prefix} (each tree node is a prefix whose KV cache all children
share --- exactly what RadixAttention-style prefix trees in serving engines provide) and
\emph{prefill-only scoring} (PRMs score prefixes, batchable across the frontier). As reasoning
models trained with RL (\cref{ch:grpo}) internalize backtracking into a single long chain of
thought, explicit search has receded for math, but the systems primitives remain central to
agentic workloads.

% ======================================================================
\section{Summary}
\begin{itemize}
  \item Scorers rank, gate and guide; they range from rules to model confidence to verifiers and
  learned reward models. Any scorer you optimize against will be gamed.
  \item Token probabilities come from teacher forcing: one forward pass over prompt plus answer,
  where logits row $i$ predicts token $i{+}1$. Confidence is not correctness.
  \item Sum log-probabilities rather than multiplying probabilities; use \texttt{log\_softmax} in
  float32 (log-sum-exp) and avoid materializing full log-softmax tensors.
  \item The average answer-token log-probability (or perplexity) is a length-normalized confidence
  score; score all candidates under the same model and prompt.
  \item Self-refinement chains draft, critique and revision calls; score gating accepts a revision
  only if it improves. Gains depend on feedback quality; intrinsic self-critique is weak on reasoning.
  \item Refinement is sequential: at a similar token budget it costs several times the latency of a
  batched vote. Order prompt content for prefix-cache hits; scoring is cheap prefill work.
\end{itemize}

% ======================================================================
\section{Exercises}

\begin{exercise}[off-by-one hunt]
Write a deliberately wrong version of \texttt{token\_logprobs} that slices logits \texttt{[1:]}
instead of \texttt{[:-1]}. Which test in Lab~16 catches it, and why would summing the wrong values
still produce plausible-looking negative numbers?
\end{exercise}

\begin{exercise}[log-sum-exp]
Prove \cref{eq:refine:lse}. Then show that \texttt{log\_softmax} computed with it never overflows for
finite inputs and that the largest output entry lies in $[-\log \vocab, 0]$.
\end{exercise}

\begin{exercise}[length bias]
Show that for any answer and any extension of it by tokens with probability $< 1$, the joint
log-probability strictly decreases. Construct an example in which the average log-probability
increases when tokens are appended. Which scorer would you use to compare a 50-token and a 500-token
solution?
\end{exercise}

\begin{exercise}[calibration]\label{ex:refine:calib}
On 100 sampled answers to problems from \texttt{math\_sample.jsonl} (or MATH-500), compute the average
log-probability and the verifier's verdict. Bin by score decile and plot accuracy per bin. Is the
model's confidence informative? Repeat with the minimum token log-probability and with the
average over boxed-answer tokens only.
\end{exercise}

\begin{exercise}[chunked log-probs]
Implement \texttt{token\_logprobs\_chunked} that computes the same result while never holding more
than $C$ rows of float32 log-softmax at once. Measure peak memory for $T = 4096$ with $C \in \{T, 512,
64\}$ on CPU using \texttt{tracemalloc} or on GPU with \texttt{torch.cuda.max\_memory\_allocated}.
\end{exercise}

\begin{exercise}[refinement harm rate]
Run self-refinement with $k=2$ on 50 problems with each of: no scorer, heuristic scorer, log-prob
scorer. For each, report (a) accuracy, (b) the fraction of initially correct drafts that end wrong,
(c) the fraction of initially wrong drafts that end correct. Which scorer has the best ratio?
\end{exercise}

\begin{exercise}[external feedback]
Replace the critique call with a verifier-based critique: if the draft's boxed answer fails a
self-check (e.g.\ substituting the answer back into an equation the model is asked to write), insert
that failure into the review. Compare against intrinsic critique.
\end{exercise}

\begin{exercise}[prefix-cache-friendly templates]
Rewrite the critique and refine templates so that the longest possible token prefix is shared across
the draft, critique and refine calls. Using the cost model of \cref{eq:refine:prefill}, compute the
prefill tokens saved for $q=150$, $a=500$, $c=120$, $k=2$, assuming perfect prefix caching.
\end{exercise}

\begin{exercise}[best-of-$n$ vs.\ refinement]
At a matched budget of $\approx$2000 decode tokens per problem, compare best-of-4 with log-prob
reranking against two rounds of refinement. Report accuracy and wall-clock latency.
\end{exercise}

% ======================================================================
\section{Lab 16: Scoring responses and score-gated refinement}\label{sec:refine:lab}

\begin{labbox}{16}{Scoring responses and score-gated refinement}
\labdir{lab16_refine}
\labsection{Goals}
Implement the scorers of this chapter, verify the teacher-forcing alignment against incremental
decoding, build the refinement loop with score gating, and quantify its cost.

\labsection{Setup}
\begin{lstlisting}[style=shell]
cd code
LAB_IMPL=starter python -m pytest -q labs/lab16_refine
\end{lstlisting}

\labsection{Tasks}
\begin{enumerate}
  \item \texttt{heuristic\_score}: format bonus plus brevity term (\cref{lst:refine:heuristic}).
  \item \texttt{answer\_token\_logprobs}: one forward pass, float32 \texttt{log\_softmax}, correct
  shift. The test compares every value against a token-by-token incremental computation.
  \item \texttt{confidence\_scores}: sum, mean, perplexity, minimum.
  \item \texttt{self\_refine}: draft, then \texttt{iterations} rounds of critique and revision with
  strict score gating; return history and the number of generation calls. The test drives it with a
  scripted fake model so acceptance and rejection are both exercised.
  \item \texttt{refine\_cost\_tokens}: the prefill/decode/sequential-call cost model of
  \cref{sec:refine:cost}.
\end{enumerate}

\labsection{Measurements}
With the Qwen3-0.6B base checkpoint, run \texttt{run.py} for 20 problems with
\texttt{--scorer none|heuristic|logprob} and \texttt{--iterations 1|2}. Record draft accuracy,
refined accuracy, acceptance rate, generated tokens and wall-clock. Expect small and inconsistent
accuracy changes with intrinsic critique, a large increase in tokens and latency per extra round, and
the log-prob scorer rejecting a sizeable fraction of revisions.

\labsection{Deliverables}
Passing tests; a table of the runs above; a short analysis of three accepted revisions that were
harmful or helpful, with the critiques that caused them.

\labsection{Stretch goals}
Implement prefix-cache-friendly templates and measure prefill savings with a prefix-caching engine
(\cref{ch:serving}); implement best-of-$n$ with log-prob reranking on top of
\texttt{sample\_n\_shared\_prefix} and compare at matched token budgets; add a calibration plot of
score versus correctness.
\end{labbox}
